\scanpage{14}
\mspage{f}
\[
\begin{aligned}
(A\times B)^\wedge&\mathop{=}\limits^{\mathrm{déf}}\underline{\mathrm{Hom}}((A\times B)^o,(\mathrm{Ens}))\simeq\underline{\mathrm{Hom}}(A^o\times B^o,(\mathrm{Ens}))\\
&\simeq\underline{\mathrm{Hom}}(A^o,\deleted{\mathrm{Ens}}\inserted{\widehat B})\simeq\underline{\mathrm{Hom}}_!(\widehat A^o,\deleted{\unsure{\mathrm{Ens}\,\widehat A}}\inserted{\widehat B})\\
&\simeq\underline{\mathrm{Hom}}(B^o,\widehat A)\simeq\underline{\mathrm{Hom}}_!(\widehat B^o,\widehat A)\\
&\simeq\underline{\mathrm{Hom}}_{!!}((\widehat A)^o\times(\widehat B)^o,\mathrm{Ens})\simeq\underline{\mathrm{Hom}}_!((\widehat A\times\widehat B)^o,\mathrm{Ens})
\end{aligned}
\]
\margin{le signe $!$ signifie : foncteurs commutant aux $\varprojlim$ quelconques ! le signe $!!$ a le sens analogue pour \uline{chacune} des variables séparément (\unsure{mais est équiv.} à un seul signe $!$ relatif à $(\widehat A\times\widehat B)^o$)}
\figure{\begin{tikzpicture}\draw (0,0)--(0,1.7);\end{tikzpicture}}
Si $X\in(A\times B)^\wedge$, le foncteur \deleted{\unsure{qui}} correspondant
\[
(\unsure{L|X},M)\longmapsto X(L,M):\underset{(\widehat A\times\widehat B)^o}{\underbrace{\widehat A^o\times\widehat B^o}}\mathbin{\unsure{\text{---}|\longrightarrow}}(\mathrm{Ens})
\]
est donné par
\[
X(L,M)=\mathrm{Hom}(L\boxtimes M,X)\deleted{=\unsure{\mathrm{Hom}}}\quad\text{où}\quad L\boxtimes M=p_B^*(L)\boxtimes p_A^*(M).
\]
\deleted{\uline{Corollaire}}\inserted{\uline{Proposition}} Le foncteur $(L,M)\longmapsto L\boxtimes M$ de $(\widehat A\times\widehat B)^o$ dans $\deleted{\unsure{(A\times B)^\wedge|(\widehat A\times\widehat B)^\wedge}}\inserted{(\widehat A\times\widehat B)^\wedge}$ est pleinement fidèle, il s'identifie au foncteur qui associe à tout \unsure{élém.} $(L,M)$ de $\widehat A\times\widehat B$ le foncteur
\deleted{\unsure{can.}\diagram{\begin{tikzcd}
(\widehat A\times\widehat B)^o\arrow[r] & \underline{\mathrm{Hom}}((\widehat A\times\widehat B)^o,
\end{tikzcd}}}
\unsure{contravariant} \deleted{\unsure{considéré qu'il représente}} sur $\widehat A\times\widehat B$ qu'il représente
\[
(L\boxtimes M)(L',M')\simeq\mathrm{Hom}_{\widehat A\times\widehat B}((L',M'),(L,M))=\mathrm{Hom}(L',L)\times\mathrm{Hom}(M',M)
\]
\margin{\unsure{Hom-}}
\[
\mathrm{Hom}(L'\boxtimes M',L^{\unsure{*}}\boxtimes M)
\]
En effet, la formule est vraie par définition de $L\boxtimes M$ si \deleted{\unsure{L'\in A}}

\deleted{Bien sûr, par fonctorialité, on a}
\deleted{\diagram{\begin{tikzcd}[column sep=large]
\mathrm{Hom}(L',L)\times\mathrm{Hom}(M',M)\arrow[r] & \mathrm{Hom}(L'\boxtimes M',L\boxtimes M)
\end{tikzcd}}}
$L'=a\in\mathrm{Ob}\,A$, $M'=b\in\mathrm{Ob}\,B$, \deleted{\unsure{et le}} D'autre part, pour $L,M$ fixés, et $L',M'$ variables, \unsure{les deux mbs}, comme fonctions de $L',M'$, transforment $\varinjlim$ en $\varprojlim$
\deleted{\unsure{ce qui est le cas}} \deleted{\unsure{et il paraît que}}
--- d'où la conclusion.

\uline{Corollaire}
\[
\underline{\mathrm{Hom}}(L'\boxtimes M',\deleted{\unsure{M}}L\boxtimes M)\simeq\underline{\mathrm{Hom}}(L',L)\boxtimes\underline{\mathrm{Hom}}(M',M)
\]
\uline{Dém.} Il suffit de prendre les $\mathrm{Hom}(a\boxtimes b,-)$, pour $(a,b)\in\mathrm{Ob}\,A\times B$
\[
\begin{array}{ll}
\Vert&\Vert\\
\mathrm{Hom}_{A\times B}(\underbrace{\deleted{a\times L'\boxtimes M'}\inserted{(a\boxtimes b)\times(L'\boxtimes M')}}_{(a\times L')\boxtimes(b\times M')},L\boxtimes M)&\mathrm{Hom}_A(a,\underline{\mathrm{Hom}}(L',L))\times\mathrm{Hom}_B(b,\underline{\mathrm{Hom}}(M',M))\\
\Vert&\Vert\\
\mathrm{Hom}_A(a\times L',L)\times\mathrm{Hom}_B(b\times M',M)&\mathrm{Hom}_A(a\times L',L)\times\mathrm{Hom}_B(b\times M',M)
\end{array}
\]
\margin{OK}
\mspage{13}
